\documentclass[conference]{IEEEtran}
\usepackage[spanish]{babel}
\usepackage[utf8]{inputenc}
\usepackage[T1]{fontenc}
\usepackage{amsmath}
\usepackage{booktabs}
\usepackage{url}

\title{Tricotomia y Comparabilidad Asintotica}
\author{\IEEEauthorblockN{Ronquillo Nunez Braulio}
\IEEEauthorblockA{Analisis y Diseno de Algoritmos}}

\begin{document}
\maketitle

\begin{abstract}
La tricotomia afirma que, para dos numeros reales, exactamente una de tres
relaciones es verdadera: menor, igual o mayor. En analisis asintotico esta
idea no siempre se conserva, porque las funciones pueden oscilar y no admitir
un orden total simple. Este reporte estudia la diferencia entre comparacion
puntual y comparacion asintotica, mostrando por que algunas funciones no son
clasificables de forma directa mediante $O$, $\Omega$ o $\Theta$.
\end{abstract}

\begin{IEEEkeywords}
tricotomia, asintotico, Big-O, limites, oscilacion.
\end{IEEEkeywords}

\section{Introduccion}
En aritmetica elemental, dos cantidades reales son comparables: si $a$ y $b$
son reales, entonces $a<b$, $a=b$ o $a>b$. Esta propiedad es natural para
numeros fijos. Sin embargo, en algoritmos se comparan funciones completas, no
solo valores individuales. Por ello, una funcion puede ser menor que otra para
algunos valores de $n$ y mayor para otros.

El uso de notaciones asintoticas como $O$ y $\Theta$ requiere observar el
comportamiento eventual de las funciones. NIST define Big-O como una medida
teorica del costo de ejecucion o memoria para un problema de tamano $n$
\cite{nist_big_o}. Esa definicion depende de constantes y de un punto $n_0$,
por lo que no se decide con una sola evaluacion.

\section{Marco Teorico}
Una relacion como $f(n)=O(g(n))$ no exige que $f(n)\leq g(n)$ para todo $n$,
sino que exista una constante $c$ tal que $f(n)\leq c g(n)$ eventualmente.
De manera similar, $f(n)=\Theta(g(n))$ exige cotas inferior y superior por
multiplos constantes de $g$ \cite{nist_theta}.

La tricotomia fallaria si se interpreta como una clasificacion total entre
clases asintoticas. Hay pares de funciones para los cuales el cociente
\begin{equation}
R(n)=\frac{f(n)}{g(n)}
\end{equation}
no tiende a cero, a una constante positiva ni a infinito. En esos casos,
la comparacion por limite no produce una de las tres conclusiones usuales.

\section{Ejemplo Oscilatorio}
Considere
\begin{equation}
f(n)=n,\qquad g(n)=n^{1+\sin n}.
\end{equation}
Entonces
\begin{equation}
\frac{f(n)}{g(n)}=n^{-\sin n}.
\end{equation}
Como $\sin n$ toma valores positivos y negativos infinitamente muchas veces,
el cociente puede hacerse muy pequeno o muy grande en distintas subsecuencias.
Por eso no hay un limite ordinario que permita clasificar el par con el
criterio simple.

MathWorld observa que incluso herramientas de limites como la regla de
L'Hospital pueden requerir cuidado cuando aparecen comportamientos
oscilatorios \cite{mathworld_lhopital}. Esta advertencia es importante:
en analisis asintotico no basta con aplicar mecanicamente una regla.

\section{Programa Implementado}
El programa simula el ejemplo anterior hasta un limite $N$. Para cada entero
$n$ calcula $f(n)$, $g(n)$ y el cociente $f(n)/g(n)$. Despues reporta el valor
minimo y maximo observados. Si los extremos se alejan de forma significativa,
se concluye que el ejemplo no se comporta como una relacion asintotica simple.

\begin{table}[htbp]
\caption{Lectura de resultados}
\centering
\begin{tabular}{ll}
\toprule
Comportamiento & Interpretacion \\
\midrule
Cociente estable & Posible $\Theta$ \\
Cociente decreciente & Posible $O$ estricta \\
Cociente creciente & Posible $\Omega$ estricta \\
Cociente oscilatorio & No hay tricotomia simple \\
\bottomrule
\end{tabular}
\end{table}

\section{Complejidad}
Si el programa evalua enteros de $2$ a $N$, realiza $N-1$ iteraciones. Cada
iteracion ejecuta operaciones aritmeticas constantes. La complejidad temporal
es $O(N)$ y la memoria adicional es $O(1)$.

\section{Conclusion}
La tricotomia de numeros reales no se traslada directamente a funciones bajo
notacion asintotica. En algoritmos, la comparacion debe demostrarse con
definiciones formales, cotas o limites bien justificados. Las funciones
oscilatorias muestran que puede no existir una clasificacion total simple.

\begin{thebibliography}{00}
\bibitem{nist_big_o} P. E. Black, ``big-O notation,'' Dictionary of Algorithms and Data Structures, NIST, 2019. [Online]. Available: \url{https://www.nist.gov/dads/HTML/bigOnotation.html}
\bibitem{nist_theta} P. E. Black, ``Theta,'' Dictionary of Algorithms and Data Structures, NIST, 2016. [Online]. Available: \url{https://www.nist.gov/dads/HTML/theta.html}
\bibitem{mathworld_lhopital} E. W. Weisstein, ``L'Hospital's Rule,'' MathWorld--A Wolfram Resource. [Online]. Available: \url{https://mathworld.wolfram.com/LHospitalsRule.html}
\bibitem{mit_algorithms} E. Demaine, J. Ku, and J. Solomon, ``6.006 Introduction to Algorithms: Lecture Notes,'' MIT OpenCourseWare, 2020. [Online]. Available: \url{https://ocw.mit.edu/courses/6-006-introduction-to-algorithms-spring-2020/pages/lecture-notes/}
\end{thebibliography}

\end{document}
